% Abstract (~280 words)
Streaming payment protocols enable continuous monetary flows between AI agents,
yet none provide flow control: when a downstream agent reaches capacity, payments
accumulate with no mechanism to reroute, buffer, or throttle---unlike data networks
where packets can be dropped or queued. We introduce \emph{Backpressure Economics}
(BPE), a cryptoeconomic mechanism that adapts the Tassiulas--Ephremides backpressure
routing algorithm from communication networks to monetary flows in multi-agent
economies, implemented as the Pura protocol.

We make four contributions. \textbf{(1)~Formal model.} We define a capacity-constrained
monetary flow network where agents declare processing capacity through
commit-reveal signals, and flow weights are set proportionally via EWMA-smoothed
capacity. \textbf{(2)~Throughput optimality.} We prove, via Lyapunov drift analysis,
that BPE achieves throughput-optimal allocation: every stabilisable demand vector
is served with bounded overflow buffers, matching the classical backpressure guarantee
in the monetary domain. \textbf{(3)~Protocol design.} We implement BPE as 32
Solidity smart contracts on Superfluid's General Distribution Agreement (GDA),
providing permissionless pool creation, concave-sqrt stake-weighted Sybil resistance,
multi-stage pipeline composition with upstream congestion propagation,
v2 extensions (NestedPool, EconomyFactory, QualityOracle, VelocityToken, UrgencyToken),
and a thermodynamic layer that maps capacity variance to a system temperature,
applies Boltzmann-weighted routing with off-chain EIP-712 aggregation,
tracks a virial ratio linking throughput to bound capital,
and triggers circuit-breaker phase transitions under sustained stress.
\textbf{(4)~Evaluation.} Simulations over 1{,}000-step horizons show BPE achieves
95.7\% allocation efficiency versus 93.5\% for round-robin under steady load,
recovers from 20\% node-kill shocks within 50 steps, and reduces buffer stall rates
from 73\% to under 9\% with overflow escrow sized at one period of peak demand.
A Sybil-cost analysis confirms that stake fragmentation yields strictly negative
net profit for all split counts $n > 1$. All contracts compile on Base Sepolia
with 319 passing unit tests across 32 test suites.
